\title{Group Sequence Policy Optimization}

\begin{abstract}
In this work, we present Group Sequence Policy Optimization (GSPO), a reinforcement learning algorithm designed to train large language models with enhanced stability, efficiency, and effectiveness. Departing from earlier methods that utilize importance ratios at the token level, GSPO introduces a sequence-wise importance ratio, facilitating clipping, reward assignment, and optimization at the sequence level. Our results show that GSPO delivers better training efficiency and superior performance over the GRPO algorithm, provides significant stabilization for Mixture-of-Experts (MoE) reinforcement learning procedures, and offers opportunities to streamline RL infrastructure design. These advances have played a pivotal role in the notable performance gains seen in recent Qwen3 models.
\end{abstract}

\section{Introduction}
\label{sec:intro}

Reinforcement learning (RL) has become a central approach in advancing the capabilities of language models \citep{o1,dpsk-r1,qwq32b,qwen3}. Employing large-scale RL techniques enables language models to address complex tasks—such as those found in competitive mathematics and programming—by engaging in more extensive and profound reasoning.

Achieving effective scaling of RL via increased computational resources critically depends on ensuring training processes remain stable and reliable. Nevertheless, leading RL algorithms such as GRPO \citep{grpo} face pronounced instability when applied to extremely large language models, frequently culminating in disastrous and unrecoverable model failures \citep{qwen3, minimax-m1}. Such training fragility poses a significant barrier to advancing the performance limits of language models through extended RL methodologies.

In this work, we demonstrate that GRPO's instability originates from a core flaw in how importance sampling weights are employed and preserved within its algorithm. This misapplication gives rise to significant high-variance noise during training, which escalates as the generated responses grow longer and is further exacerbated by the use of the clipping mechanism. Collectively, these factors lead to the eventual collapse of the model.

In response to these fundamental shortcomings, we introduce \textbf{Group Sequence Policy Optimization (GSPO)}, a novel reinforcement learning algorithm designed for the training of large language models. The principal advancement in GSPO is its rigorously formulated importance ratio, which is grounded in sequence likelihood \citep{click}, and adheres to the foundational tenets of importance sampling. Furthermore, GSPO determines normalized rewards by assessing the advantages among several responses to a single query, thereby fostering coherence between rewards assigned at the sequence level and the corresponding optimization process.

Our experimental results reveal that GSPO substantially surpasses GRPO with respect to training stability, efficiency, and overall effectiveness. Notably, GSPO inherently addresses the stability issues that commonly arise when applying RL training to large-scale Mixture-of-Experts (MoE) architectures, obviating the necessity for intricate stabilization techniques and indicating the possibility of a streamlined RL framework. These advantages have directly led to remarkable performance gains observed in recent iterations of the Qwen3 models. We anticipate that GSPO will serve as a dependable and scalable cornerstone, facilitating further progress in large-scale RL training for language models.

\section{Preliminaries}

\paragraph{Notation}
Throughout this work, we represent an autoregressive language model with parameters $\theta$ as a policy $\pi_\theta$.
The symbol $x$ signifies a query, while $\mathcal{D}$ represents the set of all queries considered. For any response $y$ provided to the query $x$, its probability according to the policy $\pi_\theta$ is expressed as $\pi_\theta (y | x)=\prod_{t=1}^{|y|} \pi_\theta (y_t | x, y_{<t} )$, where $|y|$ refers to the length of $y$ measured in tokens. The combination $(x, y)$—that is, a query paired with a corresponding response—may be evaluated by a verifier $r$, yielding a reward $r(x, y)$, which is constrained to the interval $[ 0, 1 ]$.

\paragraph{Proximal Policy Optimization (PPO)} Drawing upon data sampled from the preceding policy $\pi_{\theta_\text{old}}$, PPO \citep{ppo} maintains policy stability by limiting updates to remain close to the previous policy, accomplished through a clipping strategy. In particular, PPO optimizes the policy via the objective below (the KL regularization term is excluded here for succinctness, since it does not constitute the main subject of this work):

\begin{align}
\mathcal{J}_\text{PPO}(\theta) = \mathbb{E}_{ x \sim \mathcal{D},\, y \sim \pi_{\theta_\text{old}}( \cdot | x) }
\left[ \frac{1}{|y|} \sum_{t=1}^{|y|} 
\min \left( w_{t}(\theta) \widehat{A}_{t},  \, \mathrm{clip} \left( w_{t}(\theta), 1 - {\varepsilon}, 1 + {\varepsilon}\right) \widehat{A}_{t} \right)
\right],
\end{align}

Here, the token $y_t$'s importance ratio is given by $
w_{t}(\theta) = \frac{ \pi_{\theta} (y_{t} | x, y_{<t}) }{ \pi_{\theta_\text{old}} (y_{t} | x,y_{<t})} $, with the advantage $\widehat{A}_{t}$ of $y_t$ computed using a separate value estimator, and $\varepsilon$ denotes the interval for clipping the importance ratios.

A primary practical obstacle of PPO stems from its significant dependence on the value model. Notably, this value model is typically comparable in size to the policy model, resulting in substantial demands on both memory and computation. Moreover, the success of the algorithm is contingent upon the accuracy of its value predictions. Obtaining a trustworthy value model is intrinsically difficult; extending its capabilities to handle longer outputs and more sophisticated tasks proves to be an even more daunting problem.

\paragraph{Group Relative Policy Optimization (GRPO)} GRPO \citep{grpo} eliminates reliance on a value model by evaluating the comparative advantage of individual responses among a set of answers to an identical prompt. In particular, GRPO seeks to maximize the following objective:

\begin{align}
\label{equ:grpo}
\mathcal{J}_\text{GRPO}(\theta) = \mathbb{E}_{ x \sim \mathcal{D},\, \{y_i\}_{i=1}^G \sim \pi_{\theta_\text{old}}( \cdot | x) }
\left[ \frac{1}{G} \sum_{i=1}^{G} \frac{1}{|y_i|} \sum_{t=1}^{|y_i|} 
\min \left( w_{i,t}(\theta) \widehat{A}_{i,t}, \mathrm{clip}\left( w_{i,t}(\theta), 1 - {\varepsilon}, 1 + {\varepsilon}\right) \widehat{A}_{i,t} \right)
\right],
\end{align}

where $G$ is the number of generated responses to each query $x$ (i.e., the group size), and the importance ratio $w_{i,t}(\theta)$ and advantage $\widehat{A}_{i,t}$ of token $y_{i,t}$ are:

\begin{align}
    w_{i,t}(\theta) = \frac{ \pi_{\theta}(y_{i,t} \mid x, y_{i,<t}) }{ \pi_{\theta_\text{old}}(y_{i,t} \mid x, y_{i,<t}) }, \quad\ 
    \widehat{A}_{i,t} = \widehat{A}_{i} = \frac{ r(x, y_i) - \mathrm{mean}\left( \{ r(x, y_i) \}_{i=1}^G \right) }{ \mathrm{std}\left( \{ r(x, y_i) \}_{i=1}^G \right) },
\end{align}

respectively, where all the tokens in $y_i$ share the same advantage as $\widehat{A}_{i}$.

\section{Motivation}
\label{sec:motivation}

As models become larger, exhibit increased sparsity (such as in Mixture-of-Experts architectures), and generate longer responses, reinforcement learning demands greater rollout batch sizes to fully exploit hardware capabilities. In order to enhance sample efficiency, it is a widespread approach to divide these extensive batches of rollout data into several mini-batches for performing gradient updates. However, this strategy inevitably results in an off-policy learning environment: responses $y$ are drawn from a previous policy $\pi_{\theta_\text{old}}$ instead of from the actively optimized policy $\pi_\theta$. This circumstance also justifies why clipping techniques are integral in methods like PPO and GRPO, as they restrict the contribution of highly off-policy samples during gradient computation.

Although strategies such as clipping are intended to address the off-policy mismatch, we highlight a deeper problem within GRPO: \textit{the objective itself is fundamentally ill-posed}. This flaw is especially pronounced during the training of large-scale models on tasks that require lengthy outputs, frequently resulting in severe model degradation. The root of this ill-posedness in the GRPO objective is an incorrect usage of importance sampling weights. Importance sampling is designed to compute the expectation of a function $f$ with respect to a target distribution $\pi_\text{tar}$ by assigning appropriate weights to samples obtained from a behavior distribution $\pi_\text{beh}$:

\begin{align}
\mathbb{E}_{z \sim \pi_\text{tar}} \left[ f(z) \right] = \mathbb{E}_{z \sim \pi_\text{beh}} \left[ \frac{ \pi_\text{tar} (z) }{ \pi_\text{beh} (z)} \, f(z) \right].
\end{align}

Importantly, it is necessary to average over a large number of samples ($N \gg 1$) drawn from the behavior distribution $\pi_\text{beh}$ in order for the importance weight $\frac{ \pi_\text{tar} (z) }{ \pi_\text{beh} (z)}$ to adequately account for the discrepancy between the distributions.

Conversely, GRPO incorporates the importance weight $\frac{ \pi_{\theta} (y_{i,t} | x, y_{i,<t}) }{ \pi_{\theta_\text{old}} (y_{i,t} | x,y_{i,<t})}$ at every token timestep $t$. As this weight derives from only a single sampled token $y_{i,t}$ at each step, according to the corresponding next-token distribution $\pi_{\theta_\text{old}}( \cdot | x, y_{i, <t})$, it does not achieve its objective of correcting the distributional mismatch. Rather, it injects substantial variance into the gradient estimates, which becomes increasingly problematic in longer sequences due to the compounded noise, a situation made worse by the use of clipping. Empirical results indicate that this frequently precipitates irreversible model collapse; attempts to recover the model—such as reverting to earlier checkpoints, systematically adjusting hyperparameters like clipping thresholds, increasing sequence length, or varying RL sampling strategies—fail to restore meaningful progress once collapse has set in.

The preceding analysis indicates a foundational flaw within the architecture of GRPO. Specifically, the ineffectiveness of importance weighting at the token level highlights an essential concept: \textbf{optimization objectives must be aligned with the granularity of reward assignment}. Given that the reward is attributed to the sequence as a whole, implementing off-policy adjustments for individual tokens is likely inappropriate. Consequently, this drives us to abandon the token-level approach and consider adopting importance weighting and carrying out optimization at the \textit{sequence level} instead.

\section{Algorithm}

\subsection{GSPO: Group Sequence Policy Optimization}

Although the use of token-level importance weights $\frac{ \pi_{\theta} (y_{i,t} | x, y_{i,<t}) }{ \pi_{\theta_\text{old}} (y_{i,t} | x,y_{i,<t})}$ poses challenges in GRPO, our analysis finds that within the domain of language generation, the \textit{sequence-level} importance weight $\frac{ \pi_{\theta} (y | x) }{ \pi_{\theta_\text{old}} (y | x)}$ possesses explicit theoretical significance: it quantifies the extent to which a response $y$—sampled according to $\pi_{\theta_\text{old}} (\cdot | x)$—differs from what would be likely under $\pi_{\theta} (\cdot | x)$. This correspondence inherently matches the sequence-level reward and serves as a useful metric for evaluating the impact of clipping.

Based on this straightforward observation, we propose the \textbf{Group Sequence Policy Optimization (GSPO)} algorithm. GSPO employs the following sequence-level optimization objective:

\begin{align}
\mathcal{J}_\text{GSPO} (\theta) =
\mathbb{E}_{ x \sim \mathcal{D},\, \{y_i\}_{i=1}^G \sim \pi_{\theta_\text{old}}( \cdot | x) }
\left[ 
\frac{1}{G} \sum_{i=1}^{G}
\min \left( s_{i}(\theta)  \widehat{A}_{i}, \, \mathrm{clip} \left( s_{i}(\theta), 1 - {\varepsilon}, 1 + {\varepsilon} \right) \widehat{A}_{i} \right) 
\right],
\label{equ:gspo}
\end{align}

Here, $\mathcal{J}_\text{GSPO} (\theta)$ denotes an objective defined as an expectation over $x$ drawn from the distribution $\mathcal{D}$, and sequences $\{y_i\}_{i=1}^G$ sampled according to the policy $\pi_{\theta_\text{old}}(\cdot|x)$. Inside the expectation, an average is computed over $G$ terms, where each term involves taking the minimum between $s_i(\theta)\widehat{A}_i$ and the product of $\widehat{A}_i$ with the value of $s_i(\theta)$ limited (via $\mathrm{clip}$) to fall within $[1-\varepsilon, 1+\varepsilon]$.

where we adopt the group-based advantage estimation:

\begin{align}
\widehat{A}_{i} = \frac{r(x, y_i) - \mathrm{mean} \left( \{ r(x, y_i) \}_{i=1}^G \right) }{ \mathrm{std} \left( \{ r(x, y_i) \}_{i=1}^G \right) },
\end{align}

and define the importance ratio $s_{i}(\theta)$ based on sequence likelihood \citep{click}:

\begin{align}
s_{i}(\theta) = \left( \frac{ \pi_{\theta} (y_i | x) }{ \pi_{\theta_\text{old}} (y_i | x)} \right)^{\frac{1}{|y_i|}}
=
\exp \left( \frac{1}{|y_i|} \sum_{t=1}^{|y_i|} \log \frac{ \pi_{\theta} (y_{i,t} | x, y_{i,<t}) }{ \pi_{\theta_\text{old}} (y_{i,t} | x,y_{i,<t})} \right).
\end{align}

Consequently, GSPO implements clipping at the response level rather than on single tokens, effectively filtering out highly ``off-policy'' data from the gradient computation. This aligns the method more closely with both sequence-based reward structures and optimization objectives. It is important to highlight that length normalization is incorporated into $s_{i}(\theta)$ to mitigate variance and maintain $s_{i}(\theta)$ within a consistent numerical scale. Without such normalization, the likelihood alterations of only a few tokens might cause substantial volatility in the overall sequence importance ratio, necessitating different clipping thresholds for responses of varying lengths. Additionally, it is worth mentioning that the clipping intervals in GSPO are often of a different magnitude than those in preceding approaches, such as GRPO, primarily because of the alternative formulations of importance ratios used by these algorithms.

\subsection{Gradient Analysis}
\label{subsec:gradient}

We can derive the gradient of the GSPO objective as follows (clipping is omitted for brevity):

\begin{align}
\nabla_{\theta} \mathcal{J}_\text{GSPO} (\theta)
=&\ 
\nabla_{\theta} \mathbb{E}_{ x \sim \mathcal{D},\, \{y_i\}_{i=1}^G \sim \pi_{\theta_\text{old}}( \cdot | x) }
\left[ 
\frac{1}{G} \sum_{i=1}^{G} s_{i}(\theta) \widehat{A}_{i}
\right] \\
= &\ 
\mathbb{E}_{ x \sim \mathcal{D},\, \{y_i\}_{i=1}^G \sim \pi_{\theta_\text{old}}( \cdot | x) }
\left[ 
\frac{1}{G} \sum_{i=1}^{G}
s_{i}(\theta) \widehat{A}_{i}
\nabla_{\theta} \log s_{i}(\theta)
\right] \\
=&\ 
\mathbb{E}_{ x \sim \mathcal{D},\, \{y_i\}_{i=1}^G \sim \pi_{\theta_\text{old}}( \cdot | x) }
\left[ 
\frac{1}{G} \sum_{i=1}^{G} 
\left( \frac{ \pi_{\theta} (y_i | x) }{ \pi_{\theta_\text{old}} (y_i | x)} \right)^{\frac{1}{|y_i|}} \widehat{A}_{i} 
\cdot \frac{1}{|y_i|} \sum_{t=1}^{|y_i|} \nabla_{\theta} \log \pi_{\theta} (y_{i,t} | x, y_{i,<t}) 
\right].
\label{equ:gspo_grad}
\end{align}

For comparison, the gradient of the GRPO objective is as follows (note that $\widehat{A}_{i,t} = \widehat{A}_{i}$):

\begin{align}
\nabla_{\theta} \mathcal{J}_\text{GRPO} (\theta)
=&\
\nabla_{\theta} \mathbb{E}_{ x \sim \mathcal{D},\, \{y_i\}_{i=1}^G \sim \pi_{\theta_\text{old}}( \cdot | x) }
\left[ \frac{1}{G} \sum_{i=1}^{G} \frac{1}{|y_i|} \sum_{t=1}^{|y_i|}
w_{i,t}(\theta) \widehat{A}_{i,t}
\right] \\
=&\
\mathbb{E}_{ x \sim \mathcal{D},\, \{y_i\}_{i=1}^G \sim \pi_{\theta_\text{old}}( \cdot | x) }
\left[
\frac{1}{G} \sum_{i=1}^{G} \widehat{A}_i
\cdot \frac{1}{|y_i|} \sum_{t=1}^{|y_i|}
\frac{ \pi_{\theta} (y_{i,t} | x, y_{i,<t}) }{ \pi_{\theta_\text{old}} (y_{i,t} | x, y_{i,<t})}
\nabla_{\theta} \log \pi_{\theta} (y_{i,t} | x, y_{i,<t})
\right].
\end{align}

The key difference between GSPO and GRPO centers on \textit{the manner in which the gradients of the log likelihoods for tokens are weighted}. In GRPO, tokens are assigned weights based on their individual ``importance weight'' $\frac{ \pi_{\theta} (y_{i,t} | x, y_{i,<t}) }{ \pi_{\theta_\text{old}} (y_{i,t} | x,y_{i,<t})}$. Notably, these non-uniform weights can span ranges such as $(0, 1+\varepsilon]$ when $\widehat{A}_{i} >0$, or $[1-\varepsilon, +\infty)$ in cases where $\widehat{A}_{i} < 0$. This variability in weights is significant, and as training progresses, their cumulative effects may introduce unpredictable behavior. By contrast, GSPO addresses this by applying uniform weighting to all tokens within a response, thus removing the instability introduced by the varying weights in GRPO.

\subsection{GSPO-token: A Token-level Objective Variant}

In settings such as multi-turn RL, a more granular approach to advantage adjustment than the sequence-based method may be preferable. For this purpose, we propose a token-specific variant of GSPO, referred to as \textbf{GSPO-token}, which facilitates token-level customization of advantage:

\begin{align}
\mathcal{J}_\text{GSPO-token}(\theta)
=
\mathbb{E}_{ x \sim \mathcal{D},\, \{y_i\}_{i=1}^G \sim \pi_{\theta_\text{old}}( \cdot | x) }
\left[ 
\frac{1}{G} \sum_{i=1}^{G} \frac{1}{|y_i|} \sum_{t=1}^{|y_i|} 
\min \left( s_{i,t}(\theta) \widehat{A}_{i,t},  \, \mathrm{clip} \left( s_{i,t}(\theta), 1 - {\varepsilon}, 1 + {\varepsilon} \right) \widehat{A}_{i,t} \right) 
\right],
\label{equ:gspo-token}
\end{align}

In this formulation, the objective $\mathcal{J}_\text{GSPO-token}(\theta)$ is expressed as an expectation over input samples $x$ from the dataset $\mathcal{D}$ and $G$ trajectories $\{y_i\}$, each drawn according to $\pi_{\theta_\text{old}}(\cdot | x)$. The expression averages across both sampled sequences and over individual tokens in each sequence, resulting in the term $\frac{1}{G} \sum_{i=1}^{G} \frac{1}{|y_i|} \sum_{t=1}^{|y_i|}$. For each token at position $t$ in the $i$th trajectory, the minimum of $s_{i,t}(\theta)\,\widehat{A}_{i,t}$ and $\mathrm{clip}(s_{i,t}(\theta), 1-{\varepsilon}, 1+{\varepsilon})\,\widehat{A}_{i,t}$ is computed, ensuring stable policy updates through clipping.

where

\begin{align}
s_{i,t}(\theta) 
= \mathrm{sg} \left[ s_{i}(\theta) \right]  \times \frac{ \pi_{\theta} (y_{i,t} | x, y_{i,<t}) }{ \mathrm{sg} \left[ \pi_{\theta} (y_{i,t} | x, y_{i,<t}) \right] },
\end{align}

and $\mathrm{sg}[\cdot]$ denotes only taking the numerical value but stopping the gradient, corresponding to the \texttt{detach} operation in PyTorch. The gradient of GSPO-token can be derived as:

\begin{align}
\nabla_{\theta} \mathcal{J}_\text{GSPO-token}(\theta)
=&\ 
\nabla_{\theta} \mathbb{E}_{ x \sim \mathcal{D},\, \{y_i\}_{i=1}^G \sim \pi_{\theta_\text{old}}( \cdot | x) }
\left[
\frac{1}{G} \sum_{i=1}^{G}
\frac{1}{|y_i|} \sum_{t=1}^{|y_i|}
s_{i,t}(\theta) \widehat{A}_{i,t}
\right] \\
=&\
\mathbb{E}_{ x \sim \mathcal{D},\, \{y_i\}_{i=1}^G \sim \pi_{\theta_\text{old}}( \cdot | x) }
\left[
\frac{1}{G} \sum_{i=1}^{G} s_i (\theta) \cdot \frac{1}{|y_i|} \sum_{t=1}^{|y_i|}
\widehat{A}_{i,t} \frac{\nabla_{\theta} \pi_{\theta} (y_{i,t} | x, y_{i,<t}) }{\pi_{\theta} (y_{i,t} | x, y_{i,<t})}
\right] \\
=&\
\mathbb{E}_{ x \sim \mathcal{D},\, \{y_i\}_{i=1}^G \sim \pi_{\theta_\text{old}}( \cdot | x) }
\left[
\frac{1}{G} \sum_{i=1}^{G} \left( \frac{ \pi_{\theta} (y_i | x) }{ \pi_{\theta_\text{old}} (y_i | x)} \right)^{1 / |y_i|}
\cdot \frac{1}{|y_i|} \sum_{t=1}^{|y_i|}
\widehat{A}_{i,t} \nabla_{\theta} \log \pi_{\theta} (y_{i,t} | x, y_{i,<t})
\right].
\label{equ:gspo-token_grad}
\end{align}

Observe that the expression $\frac{ \pi_{\theta} (y_{i,t} | x, y_{i,<t}) }{ \mathrm{sg} \left[ \pi_{\theta} (y_{i,t} | x, y_{i,<t}) \right] }$ evaluates to 1; as a result, $s_{i,t}(\theta)$ and $s_{i}(\theta)$ are equal in terms of their numerical value.
When we examine Equation~\eqref{equ:gspo} alongside~\eqref{equ:gspo-token}, and Equation~\eqref{equ:gspo_grad} with~\eqref{equ:gspo-token_grad}, it becomes evident that GSPO-token and GSPO coincide exactly in terms of their optimization objectives, clipping mechanisms, and formal gradients whenever we set the advantages of every token in response $y_i$ uniformly (specifically, $\widehat{A}_{i,t} = \widehat{A}_{i}$ for all $t$). However, GSPO-token offers the added benefit of enabling token-wise customization of advantages, thus providing greater flexibility.

\section{Experiments and Discussion}

\subsection{Empirical Results}

We conduct experiments utilizing a cold-start model derived from fine-tuning Qwen3-30B-A3B-Base, and present both the training reward trajectories and model performance results across the AIME'24 (mean Pass@1 from 32 samples), LiveCodeBench (202410-202502, mean Pass@1 from 8 samples), and CodeForces (Elo Rating) evaluation suites. For RL training, each rollout data batch is subdivided into four mini-batches for gradient update steps. The GSPO approach implements left and right clipping thresholds in Equation~\eqref{equ:gspo} at 3e-4 and 4e-4, accordingly. For comparison, we utilize GRPO as our reference baseline, employing left and right clipping limits in Equation~\eqref{equ:grpo} set to 0.2 and 0.27, values selected for optimal comparability. It is important to highlight that GRPO requires the application of the Routing Replay training mechanism to ensure standard convergence in MoE RL, a detail further examined in \S~\ref{subsec:routing_replay}; in contrast, \textit{GSPO eliminates the necessity for this procedure}.

As illustrated in Figure~\ref{fig:results}, GSPO enables consistently stable training dynamics. Notably, \textbf{GSPO consistently enhances performance as training compute is scaled up, the query set is periodically refreshed, and the sequence length is increased}. In addition, GSPO outperforms GRPO in terms of training efficiency, as evidenced by higher training accuracy and superior benchmark outcomes while utilizing equivalent computational resources and numbers of queries. Ultimately, the application of GSPO to RL training for the latest Qwen3 models provides compelling evidence of its effectiveness, affirming GSPO's capacity to realize the benefits of RL scaling in large language model development.

\subsection{Curious Observation on Clipping Fractions}

One of the principal differences between GSPO and GRPO lies in GSPO's strategy of clipping responses at the sequence level rather than acting on each token individually. As illustrated in Figure~\ref{fig:clipping}, this methodology results in a discrepancy of approximately two orders of magnitude in the proportions of clipped tokens between the two approaches (and varying the clipping thresholds does not bridge this substantial gap). Despite the fact that GSPO clips a much higher proportion of tokens—thereby utilizing a smaller subset of data for model updates or gradient estimation—it paradoxically outperforms GRPO in terms of training efficiency. This unexpected result, where substantially increased clipping leads to better efficiency, implies that token-wise gradient signals in GRPO tend to be inherently high in variance and less effective for learning from samples. By operating at the sequence level, GSPO produces learning signals that are both more stable and more conducive to efficient optimization.

\subsection{Benefit of GSPO for MoE Training}
\label{subsec:routing_replay}

\paragraph{Background}
Unlike standard RL training on dense models, the inherently sparse activations in MoE architectures present distinct challenges regarding stability. Notably, our investigations revealed that employing the GRPO algorithm often leads to pronounced \textit{expert-activation volatility}, which can undermine the convergence of RL optimization in MoE models. Specifically, after one or more gradient updates, the subset of experts selected for the identical response may shift appreciably. For instance, in the 48-layer Qwen3-30B-A3B-Base, we observe that following each RL step, about 10\% of the experts newly activated by the updated policy $\pi_{\theta}$ differ from those chosen by the preceding policy $\pi_{\theta_\text{old}}$ for the same rollout input. This volatility becomes increasingly severe as MoE models scale in depth, causing significant instability in the token-level importance weights $w_{i,t}(\theta) = \frac{ \pi_{\theta} (y_{i,t} | x, y_{i,<t}) }{ \pi_{\theta_\text{old}} (y_{i,t} | x,y_{i,<t})}$, rendering them unreliable—as examined in \S~\ref{sec:motivation} and~\ref{subsec:gradient}. This instability, in turn, impedes the regular convergence behavior of RL training.

\paragraph{Our Previous Approach} In our earlier work, we addressed this issue by utilizing the \textbf{Routing Replay} training methodology. Concretely, this involved storing the active experts determined by $\pi_{\theta_\text{old}}$ and subsequently “replaying” these identical routing decisions during calculations under $\pi_{\theta}$ for the computation of the importance ratios $w_{i,t}(\theta) = \frac{ \pi_{\theta} (y_{i,t} | x, y_{i,<t}) }{ \pi_{\theta_\text{old}} (y_{i,t} | x,y_{i,<t})}$. Through this approach, both $\pi_{\theta} (y_{i,t} | x, y_{i,<t})$ and $\pi_{\theta_\text{old}} (y_{i,t} | x,y_{i,<t})$ utilize an identical network configuration per token $y_{i,t}$, which helps maintain consistency in token-wise importance ratios and guarantees that optimization proceeds over a stable set of active network paths throughout gradient steps. As illustrated in Figure~\ref{fig:routing_replay}, Routing Replay proves to be a critical component in ensuring smooth GRPO training convergence in MoE models.

\paragraph{Benefit of GSPO} While Routing Replay facilitates effective convergence during GRPO training of MoE models, it does so at the expense of introducing additional memory and communication costs, and can constrain the true capacity available to the MoE model. In comparison, as illustrated in Figure~\ref{fig:results}, GSPO removes any reliance on Routing Replay and permits standard computation of importance ratios $s_i(\theta)$, thereby supporting stable optimization and regular convergence. The central insight underpinning GSPO is its exclusive consideration of sequence likelihood (i.e., $\pi_{\theta} (y_i | x)$), with minimal sensitivity to token-level likelihoods (i.e., $\pi_{\theta} (y_{i,t} | x, y_{i,<t})$). As the MoE model continually preserves its underlying language modeling abilities, the associated sequence likelihood exhibits minimal variability. Thus, GSPO directly addresses the problem of expert-activation instability characteristic of MoE models, dispensing with the need for elaborate solutions such as Routing Replay. This approach not only enhances the robustness and efficiency of training but also empowers the model to make full use of its representational capacity.

\subsection{Benefit of GSPO for RL Infrastructure}

Due to differences in numerical precision between training frameworks such as Megatron and inference systems like SGLang or vLLM, it is common practice to re-evaluate the likelihoods of sampled outputs under the previous policy $\pi_{\theta_\text{old}}$ using the training engine. Nevertheless, GSPO focuses on optimizing sequence-level likelihoods instead of token-level ones; the former are generally less sensitive to precision inconsistencies. As a result, GSPO enables the use of likelihoods directly obtained from inference engines during optimization, eliminating the necessity to perform recomputation via the training engine. This approach proves particularly advantageous for tasks involving partial rollouts, multi-turn reinforcement learning, and frameworks that separate training and inference processes.

\section{Conclusion}

We introduce Group Sequence Policy Optimization (GSPO), a novel reinforcement learning approach tailored for optimizing large language models. Rooted in the framework of importance sampling, GSPO establishes importance ratios via sequence likelihoods and implements sequence-wise procedures for clipping, reward assignment, and policy improvement. When contrasted with GRPO, GSPO offers marked advancements in training stability, computational efficiency, and overall performance. Its advantages are especially significant in large-scale reinforcement learning applications involving MoE models, and it underpins the remarkable progress observed in the latest Qwen3 systems. As GSPO provides a robust and extensible algorithmic platform, we aim to further scale reinforcement learning and anticipate substantial breakthroughs in artificial intelligence as a result.

\end{document}